\documentclass[10pt,twocolumn]{article}
\usepackage[margin=0.75in,columnsep=0.3in]{geometry}
\usepackage{amsmath,amssymb}
\usepackage{graphicx}
\usepackage{booktabs}
\usepackage{tabularx}
\usepackage[hidelinks]{hyperref}

\title{\vspace{-1em}Numerical Analysis of the SIR Epidemic Model:\\
Solver Accuracy, Parameter Recovery, and Robustness to Noise}
\author{
  Saif Uz Zaman \quad Sayaad Muzahid Masfi \quad Ibtida bin Ahmed \quad Sakif Naieb Raiyan \quad Aurchi Chowdhury \\
  \small CSE 402: Numerical Analysis, Simulation \& Modeling --- Section B, Group 2
}
\date{}

\begin{document}
\maketitle

\begin{abstract}
Parameter estimation for the SIR epidemic model is typically studied with the ODE solver treated
as a black box. We make solver choice the central experimental variable, implementing Forward
Euler, Heun's method, classical RK4, adaptive Dormand--Prince RK45, and two implicit methods from
scratch, and asking how solver accuracy propagates into least-squares-recovered transmission and
recovery rates ($\beta$, $\gamma$). Using a semi-analytic gold-standard SIR solution (derived from
the model's exact phase-plane reduction) as ground truth, we show that estimation bias inherits the
fitting solver's convergence order, and we locate a crossover noise level $\sigma^*$ at which this
solver-induced bias is exactly matched by noise-induced estimator variance: below $\sigma^*$,
solver choice dominates the error budget; above it, it is statistically irrelevant. The pipeline
draws on all four CSE 402 lab assignments and the optimization home assignment -- root finding,
LU and QR factorization, eigenvalues by the power method, random number generation and Monte
Carlo simulation, optimization, least squares, interpolation, and quadrature. We map each method to where it is used and document how
every random number is generated. All code, experiments, and figures in this report are regenerated from a single
reproducible pipeline (\texttt{make all}).
\end{abstract}

\section{Introduction}
The SIR (Susceptible--Infectious--Recovered) model is the canonical compartmental description of an
epidemic, and fitting it to case data is a standard tool for estimating a pathogen's transmission
rate $\beta$ and recovery rate $\gamma$, and hence its basic reproduction number $R_0 = \beta/\gamma$.
Capaldi et al.~\cite{capaldi2012} study parameter estimation and uncertainty quantification for this
model in detail: least-squares fitting, sensitivity of the estimates to $\beta$--$\gamma$ correlation,
and the effect of sampling frequency. Like most work in this area, they treat the numerical
integration of the ODE system as a solved problem -- a black box that produces ``the'' model
trajectory, with no further scrutiny of \emph{how} it is produced.

We ask the question that black box leaves closed: \emph{does the choice and accuracy of the ODE
solver materially affect the parameters recovered from it}, once realistic imperfections --
observation noise, sparse sampling, an incompletely observed epidemic -- are introduced? We answer
it by implementing every solver ourselves (Forward Euler, Heun, classical RK4, adaptive
Dormand--Prince RK45, and two implicit methods), building an estimation pipeline around them with
hand-written optimizers, and running a controlled factorial experiment that isolates solver
truncation error from every other source of error in the pipeline.

Our central methodological device is a \emph{semi-analytic gold standard}: the SIR system admits an
exact reduction to a single quadrature (Section~\ref{sec:methods}), so we can measure every solver's
error against the true solution to machine precision, rather than against ``a finer numerical
solve'' -- which would be circular. This lets us make a quantitative claim that a black-box
treatment cannot: there is a \emph{crossover noise level} $\sigma^*$ at which solver-induced
estimation bias equals noise-induced estimation variance. Below $\sigma^*$, solver accuracy is the
dominant source of error in the recovered parameters, and using Forward Euler at a coarse step is
scientifically indefensible. Above $\sigma^*$, sampling noise dominates and a cheaper, less accurate
solver is statistically indistinguishable from a more expensive one.


\section{Background}
\subsection{The SIR model}
The classical SIR model partitions a closed population of size $N$ into susceptible ($S$),
infectious ($I$), and recovered ($R$) compartments:
\begin{align}
  \frac{dS}{dt} &= -\frac{\beta S I}{N}, &
  \frac{dI}{dt} &= \frac{\beta S I}{N} - \gamma I, &
  \frac{dR}{dt} &= \gamma I. \label{eq:sir}
\end{align}
The basic reproduction number is $R_0 = \beta/\gamma$: an epidemic grows while $R_0 \cdot S(t)/N > 1$
and declines once susceptible depletion pushes it below 1.

\subsection{Base paper and our extension}
Capaldi et al.~\cite{capaldi2012} fit \eqref{eq:sir} to synthetic and real incidence data via
least squares, quantify the uncertainty in the resulting $(\hat\beta,\hat\gamma)$ via an asymptotic
covariance estimate, and study how estimation quality degrades with sparser sampling. Their solver
is used only as a means to an end. We take their estimation machinery as a starting point but
restructure the experiment so that \emph{solver identity and step size} are first-class independent
variables, crossed with noise level, sampling interval, observation window, and the underlying
$(\beta,\gamma)$ regime.

\subsection{Related numerical-methods reference}
Prodanov~\cite{prodanov2021} gives an analytical treatment of the SIR system's phase-plane
reduction that we build on directly for our gold-standard solution (Section~\ref{sec:methods}).


\section{Methods}
\label{sec:methods}
\subsection{Semi-analytic gold standard}
Dividing $dS/dt$ by $dR/dt$ in \eqref{eq:sir} gives $dS/dR = -(\beta/\gamma N) S$, so
\begin{equation}
  S(t) = S_0 \exp\!\Big(-\tfrac{\beta}{\gamma N}\big(R(t)-R_0\big)\Big), \label{eq:phase}
\end{equation}
an exact, solver-free relation between $S$ and $R$ (we write $R_0$ here for the initial value of the
$R$ compartment, not to be confused with the reproduction number). Substituting \eqref{eq:phase} into
$dR/dt = \gamma I = \gamma(N-R-S)$ gives a \emph{scalar} separable ODE, so
\begin{equation}
  t(R) = \int_0^{R} \frac{dr}{\gamma\big(N - r - S_0 e^{-\frac{\beta}{\gamma N} r}\big)}. \label{eq:tofr}
\end{equation}
We evaluate \eqref{eq:tofr} by adaptive Simpson quadrature and invert it by a safeguarded
Newton method (the standard \texttt{rtsafe} scheme) to obtain $R(t)$, then recover $S$ from
\eqref{eq:phase} and $I = N-S-R$. This reference trajectory is accurate to machine precision using
no ODE solver at all, and every accuracy claim in this paper is measured against it. We additionally
use \eqref{eq:phase} to define the \emph{phase invariant}
$Q(t) = \ln S(t) + \frac{\beta}{\gamma N} R(t)$, constant along any true trajectory; unlike the
trivial mass-conservation check $S+I+R=N$ (which every explicit Runge--Kutta method satisfies to
roundoff at \emph{any} step size, since the right-hand side of \eqref{eq:sir} sums to zero
identically), drift in $Q(t)$ scales with the solver's true convergence order and is a genuine
accuracy diagnostic.

\subsection{Solvers}
We implement Forward Euler (order 1), Heun's method (order 2), classical RK4 (order 4), adaptive
Dormand--Prince RK45 with a PI step controller (order 5(4), used as the ``truth'' solver for model
variants with no closed form), and two implicit methods -- backward Euler and the trapezoidal rule
-- whose nonlinear stage equations are solved by Newton's method with the linear system at each
iteration factored by our own partial-pivoting LU decomposition.

\subsection{Estimation}
Given noisy observations $\mathrm{obs}_k$ at times $t_k$, we minimize the weighted least-squares cost
\begin{equation}
  J(\theta) = \sum_k w_k \big(\mathrm{obs}_k - h(y(t_k;\theta))\big)^2 \label{eq:cost}
\end{equation}
by grid search, golden-section coordinate descent, Nelder--Mead, Gauss--Newton, and
Levenberg--Marquardt, all implemented from scratch. The Jacobian required by the last two methods is
obtained by integrating the \emph{forward sensitivity equations}
$\dot s_j = J_y f \, s_j + \partial f/\partial\theta_j$ alongside the state, using the same solver
under test -- so solver error enters the Jacobian exactly as it enters the trajectory, which is the
mechanism behind our central claim, rather than by finite differences.

\subsection{Uncertainty quantification}
We compute four independent confidence regions for $(\hat\beta,\hat\gamma)$: an asymptotic ellipse
from $\mathrm{Cov}(\hat\theta) \approx \hat\sigma^2 (J^\top J)^{-1}$; a residual bootstrap; a Monte
Carlo ensemble over fresh noise realizations from the known ground truth (the only route that
measures true bias); and an adaptive Metropolis--Hastings sampler with four chains, reporting the
Gelman--Rubin $\hat R$ and effective sample size.

\subsection{The crossover experiment}
For a fixed $(\beta,\gamma)$, sampling interval, and fitting solver/step size $h$, we generate
noise-free synthetic data from the gold standard and fit with that solver: any deviation from the
true parameters is attributable purely to the solver's truncation error, giving a deterministic
$\mathrm{bias}(h)$. We then add noise at level $\sigma$ and refit over many replicates, giving
$\mathrm{std}(\sigma, h)$. The crossover $\sigma^*(h)$ is where $\mathrm{std}(\sigma^*, h) =
|\mathrm{bias}(h)|$.


\section{CSE 402 Methods Used}
\label{sec:course}
Table~\ref{tab:course} maps each topic of the CSE 402 lab outline to the method we implemented
for it and to where that method is used. The project covers all four lab assignments and the
optimization home assignment, and so all three course outcomes: nonlinear equations and linear
systems (CO1), interpolation, regression, and numerical integration (CO2), and simulation with
random number generation and Monte Carlo methods (CO3). Unless marked $^\dagger$, every method in
the table is our own code in \texttt{core/sirlab}, which never imports SciPy (enforced by
\texttt{tests/test\_hygiene.py}); SciPy appears only in tests, as an independent cross-check.

\begin{table*}[tp]
  \centering
  \footnotesize
  \caption{CSE 402 syllabus topics and where the project uses them. Code paths are relative to
    \texttt{core/sirlab/}; \texttt{E}-numbered scripts are in \texttt{experiments/}.
    $^\dagger$\,Computed with a NumPy routine rather than our own code.}
  \label{tab:course}
  \begin{tabularx}{\textwidth}{@{}>{\raggedright\arraybackslash}p{2.7cm}
      >{\raggedright\arraybackslash}X >{\raggedright\arraybackslash}p{3.5cm}@{}}
    \toprule
    Syllabus topic & What we implemented, and where it is used & Code \\
    \midrule
    \multicolumn{3}{@{}l}{\textbf{Lab 1: approximations, errors, and root finding (CO1)}} \\
    Truncation error & Global error vs.\ the gold standard; observed order $\hat p$ (E01); bias in
      $\hat\beta \propto h^p$ (E07) & \texttt{diagnostics.py} \\
    Round-off error & $S+I+R=N$ holds to round-off for every explicit RK method at any $h$ (E02);
      adaptive Simpson stops refining at the round-off floor
      & \texttt{diagnostics.py}, \texttt{quadrature/adaptive.py} \\
    Bisection, Newton--Raphson & Safeguarded Newton (\texttt{rtsafe}): Newton steps inside the
      bracket, bisection otherwise; inverts $t(R)$ and solves the final-size equation
      & \texttt{roots/newton.py} \\
    Newton--Raphson for systems & $(I - hJ_f)\,\delta = -G(y)$ on each backward-Euler or
      trapezoidal stage & \texttt{solvers/implicit.py} \\
    Plotting & All figures generated from the experiment JSON artifacts
      & \texttt{scripts/make\_figures.py} \\
    \midrule
    \multicolumn{3}{@{}l}{\textbf{Lab 2: systems of linear equations and eigenvalues (CO1)}} \\
    Gauss elimination, LU decomposition & Doolittle LU with partial pivoting, forward/back
      substitution; every implicit Newton step, the GN/LM normal equations, every inverse-power step
      & \texttt{linalg/lu.py} \\
    Matrix inverse, conditioning & $(J^\top J)^{-1}$ for the asymptotic covariance via LU solves;
      1-norm $\kappa(J^\top J)$ (E05); 2-norm $\kappa(J)$ from the eigenvalues of $J^\top J$ (E10)
      & \texttt{linalg/lu.py}, \texttt{linalg/eigen.py} \\
    Power method & Largest-entry normalization; $h\,\rho(J_y f)$ for stability (E03), with a 2-D
      projection when the SIR Jacobian's eigenvalues are complex (Section~\ref{sec:eigen})
      & \texttt{linalg/eigen.py}, \texttt{diagnostics.py} \\
    Inverse power method & Smallest eigenvalue by LU solves, factored once; $\lambda_{\min}(J^\top J)$
      for $\kappa_2(J)$ (E10) & \texttt{linalg/eigen.py} \\
    Hotelling deflation & All eigenpairs of a symmetric matrix: $\mathrm{Cov}(\hat\theta)$ for the
      ellipse axes (E11, E14); $J^\top J$ for the cost valley (E05)
      & \texttt{linalg/eigen.py}, \texttt{uq/asymptotic.py} \\
    \midrule
    \multicolumn{3}{@{}l}{\textbf{Lab 3: random number generation and Monte Carlo methods (CO3)}} \\
    Random number generation & PCG64 uniform stream; exact normal, uniform, integer, and gamma
      variates; per-replicate seeds (Section~\ref{sec:rng})$^\dagger$
      & \texttt{observe.py}, E05--E11, E13 \\
    Monte Carlo simulation & Replicated synthetic datasets from the known truth, each refitted:
      bias, variance, RMSE (E07--E10, E13); UQ route 3 (E11)
      & \texttt{uq/montecarlo.py}, E07--E11, E13 \\
    Bootstrap resampling & Residual bootstrap of $(\hat\beta, \hat\gamma)$ (E11); percentile band
      on the RMSE curve (E08) & \texttt{uq/bootstrap.py}, \texttt{E08\_noise.py} \\
    Metropolis--Hastings & Adaptive random-walk Metropolis in $\ln\theta$ plus a Gibbs step for
      $\sigma^2$; 4 chains, acceptance rate, $\hat R$~\cite{gelman1992}, ESS (E11)
      & \texttt{uq/mcmc.py} \\
    \midrule
    \multicolumn{3}{@{}l}{\textbf{Home Assignment 1: QR and optimization (CO1)}} \\
    QR & Householder QR step $R\,\delta = -Q^\top r$ avoids the $\kappa(J)^2$ of the normal
      equations; unit-tested, but reported fits use LU. Not the QR eigenvalue iteration
      & \texttt{linalg/qr.py}, \texttt{estimate/gaussnewton.py} \\
    Golden-section search & 1-D minimizer in coordinate descent (E06) and the profile likelihoods
      (E11) & \texttt{estimate/golden.py}, \texttt{uq/profile.py} \\
    Newton's method & Gauss--Newton: Newton with Hessian $\approx J^\top J$, backtracking line
      search (E06) & \texttt{estimate/gaussnewton.py} \\
    Gradient methods & Levenberg--Marquardt spans Gauss--Newton ($\lambda \to 0$) and scaled
      gradient descent ($\lambda \to \infty$); the estimator in E07--E11, E13, E14
      & \texttt{estimate/lm.py} \\
    Constrained optimization & $\beta, \gamma > 0$ by optimizing $u = \ln\theta$; box bounds in
      golden-section search & \texttt{estimate/lm.py}, \texttt{estimate/golden.py} \\
    Beyond the list & Nelder--Mead: derivative-free baseline (E06) and 3-parameter fit (E10);
      $J(\beta, \gamma)$ on a 2-D grid (E05)
      & \texttt{estimate/neldermead.py}, \texttt{E05\_landscape.py} \\
    \midrule
    \multicolumn{3}{@{}l}{\textbf{Lab 4: curve fitting, interpolation, and numerical integration (CO2)}} \\
    Least-squares regression & Nonlinear least squares \eqref{eq:cost} in every fit; linear least
      squares on $(\ln h, \ln e)$ gives $\hat p$ (E01) & \texttt{estimate/}, \texttt{diagnostics.py} \\
    Linear interpolation & Locates the crossover $\sigma^*$ between grid points (E13)
      & \texttt{E13\_crossover.py} \\
    Spline interpolation & Piecewise cubic Hermite dense output evaluates trajectories and
      sensitivities at the observation times & \texttt{solvers/base.py} \\
    Trapezoidal and Simpson's rules & Adaptive Simpson for the gold-standard integral
      \eqref{eq:tofr}. Composite and uneven-grid Simpson (the integral of the 3-point Lagrange
      quadratic) and the trapezoidal rule serve the incidence operator: unit-tested, but unused in
      the reported experiments & \texttt{quadrature/adaptive.py}, \texttt{quadrature/simpson.py} \\
    Richardson, Romberg & Each accepted Simpson panel adds $(S_2 - S_1)/15$ (one- vs.\ two-panel
      estimates): one Romberg step & \texttt{quadrature/adaptive.py} \\
    Numerical differentiation & Central-difference node slopes for the Hermite
      interpolant$^\dagger$; complex-step checks of the analytic Jacobians
      & \texttt{solvers/base.py}, \texttt{tests/test\_models.py} \\
    \midrule
    \multicolumn{3}{@{}l}{\textbf{Beyond the lab outline: ODE solvers (course textbook~\cite{chapra2015}, Part Seven)}} \\
    Runge--Kutta methods & Euler, Heun, and classical RK4, with their stability frontier (E03);
      adaptive Dormand--Prince RK45 with PI step control
      & \texttt{solvers/explicit.py}, \texttt{solvers/rk45.py} \\
    Implicit methods & Backward Euler, trapezoidal rule (A-stable), in the convergence study (E01)
      & \texttt{solvers/implicit.py} \\
    \bottomrule
  \end{tabularx}
\end{table*}

Two further points complete the picture. First, the ODE solvers are not a CSE 402 lab topic: we
took them from Part Seven of the course textbook~\cite{chapra2015}, and they are the project's
experimental variable rather than its apparatus. Second, a few syllabus methods have no role in
this problem and are not used: false position, Bairstow's method, Gauss--Jordan elimination, the
QR eigenvalue iteration, and Newton's divided-difference interpolation.


\section{Results}
\label{sec:results}
\emph{Numbers in this section are from the \texttt{quick} experiment profile (Section~\ref{sec:appendix-profiles}); the \texttt{full} profile matches the factorial sizes in our project plan and is intended for an offline run.}

\subsection{RQ1: solver behavior}
Figure~\ref{fig:e01} shows the fitted convergence order for each solver against the gold standard.
Observed orders match theory closely: Forward Euler $\hat p=1.03$ (theory 1), Heun $\hat p=1.94$
(theory 2), classical RK4 $\hat p=3.94$ (theory 4), backward Euler $\hat p=0.97$ (theory 1), and
trapezoidal $\hat p=2.00$ (theory 2).

\begin{figure}[t]
  \centering
  \includegraphics[width=\linewidth]{../figures/e01_convergence.pdf}
  \caption{Observed convergence order matches theory (E01, $R_0=3$ regime).}
  \label{fig:e01}
\end{figure}

Figure~\ref{fig:e02} demonstrates the deliberate finding of Section~\ref{sec:methods}: mass error
stays at roundoff magnitude regardless of solver or step size, while the phase-invariant drift
correctly separates solver quality.

\begin{figure}[t]
  \centering
  \includegraphics[width=\linewidth]{../figures/e02_invariants.pdf}
  \caption{Mass conservation is uninformative; phase-invariant drift is not (E02).}
  \label{fig:e02}
\end{figure}

\subsection{RQ2: solver-in-the-loop parameter recovery}
Figure~\ref{fig:e07} is the mechanism behind our central claim: with noise-free synthetic data
(so any deviation from truth is attributable purely to the fitting solver), bias in the recovered
$\hat\beta$ shrinks with $h$ at each solver's own convergence order. At $(\beta,\gamma)=(0.3,0.1)$,
reducing $h$ from $0.25$ to $0.05$ (a factor of 5) shrinks Euler's bias by a factor of $\approx 5.0$
($4.48\times10^{-3} \to 8.93\times10^{-4}$), Heun's by $\approx 24.2$
($5.88\times10^{-5} \to 2.43\times10^{-6}$), and RK4's by $\approx 606$
($7.17\times10^{-9} \to 1.18\times10^{-11}$) -- close to the theoretical $5^1$, $5^2 = 25$, and
$5^4 = 625$.

\begin{figure}[t]
  \centering
  \includegraphics[width=\linewidth]{../figures/e07_recovery_bias.pdf}
  \caption{Solver-induced bias in $\hat\beta$ scales with the solver's convergence order (E07).}
  \label{fig:e07}
\end{figure}

\subsection{RQ3: noise, sampling, and uncertainty}
Figure~\ref{fig:e08} shows RMSE$(\hat\beta)$ growing with noise level $\sigma$, with bootstrap
confidence bands on the RMSE curve itself; at our reference solver/step size, RMSE grows from
$\approx 1.2\times10^{-11}$ (noise-free) to $5.6\times10^{-4}$ at $\sigma=0.02$ and
$2.9\times10^{-3}$ at $\sigma=0.10$.

\begin{figure}[t]
  \centering
  \includegraphics[width=\linewidth]{../figures/e08_noise_degradation.pdf}
  \caption{RMSE$(\hat\beta)$ vs.\ noise level, with bootstrap CI (E08).}
  \label{fig:e08}
\end{figure}

Figure~\ref{fig:e09} shows RMSE$(\hat\beta)$ as the observation window shrinks from 6 to 1.5 times
the peak time (E09). With dense sampling, truncation costs little: RMSE rises by 34\% at
$dt = 0.5$ and 22\% at $dt = 1$. At $dt = 2$ it doubles ($1.5\times10^{-3} \to 3.0\times10^{-3}$),
so coarse sampling and a short window compound. Every window in this grid still contains the peak;
cutting the observations off \emph{before} the peak is a far stronger effect
(Figure~\ref{fig:e11}). Separately (E10), treating $I_0$ as an unknown third parameter rather than
a known constant raises the condition number of the Jacobian from $\kappa_2 \approx 2.2$--$2.6$ to
$\approx 211$ at $I_0 = 1$ and $\approx 1828$ at $I_0 = 20$. It inflates the variance of
$\hat\beta$ by a factor of 19 at $I_0 = 1$ and 13 at $I_0 = 5$, but only 3.5 at $I_0 = 20$
(10 replicates per cell), so the growth in condition number does not carry over to the variance.

\begin{figure}[t]
  \centering
  \includegraphics[width=\linewidth]{../figures/e09_sampling_window.pdf}
  \caption{RMSE$(\hat\beta)$ vs.\ observation window (as a fraction of 6 times the peak time) and
    sampling interval (E09).}
  \label{fig:e09}
\end{figure}

\subsection{Uncertainty quantification}
Figure~\ref{fig:e11} overlays all four UQ routes and compares profile likelihoods on the full vs.\
a window truncated at 0.3 times the peak time. The truncated profile is nearly flat: it varies by
10\% across the $\beta$ grid, against a 65-fold rise for the full window. This is the signature of
non-identifiability before the peak.

\begin{figure}[t]
  \centering
  \includegraphics[width=\linewidth]{../figures/e11_uncertainty.pdf}
  \caption{Four UQ routes agree closely (left); a truncated window flattens the profile likelihood, the visual signature of non-identifiability (right) (E11).}
  \label{fig:e11}
\end{figure}

\subsection{The crossover \texorpdfstring{$\sigma^*$}{σ*}}
Figure~\ref{fig:e13} overlays the deterministic solver bias in $\hat\beta$ against the
noise-induced standard deviation for Forward Euler at $h = 0.1$; they cross at $\sigma^* = 0.058$
(Section~\ref{sec:methods}). At $h = 0.25$ the bias is 2.5 times larger ($4.48\times10^{-3}$, the
same value E07 finds) and $\sigma^* = 0.156$, so $\sigma^*$ grows roughly in proportion to $h$: a
first-order method's bias scales as $h$, while the noise-induced spread grows with $\sigma$. At
$h = 0.5$ the spread is still below the bias at the largest noise level tested ($\sigma = 0.2$), so
there $\sigma^* > 0.2$. For $\hat\gamma$ the crossover comes earlier: $\sigma^* = 0.029$, $0.054$,
and $0.128$ at $h = 0.1$, $0.25$, and $0.5$.

\begin{figure}[t]
  \centering
  \includegraphics[width=\linewidth]{../figures/e13_crossover.pdf}
  \caption{Crossover for Forward Euler at $h = 0.1$: below $\sigma^* = 0.058$ solver bias dominates
    the error in $\hat\beta$; above it, noise dominates (E13).}
  \label{fig:e13}
\end{figure}

\subsection{Extensions: model variants and real data}
The SEIR, SIRS, and SIR-with-vaccination variants (E12) preserve the RQ1 solver ordering
(Euler~$<$~Heun~$<$~RK4 in accuracy at fixed $h$), indicating the finding is not an artifact of the
plain SIR system's particular structure. Applying the same pipeline to the 1666 Eyam plague outbreak
(E14) recovers an $\hat R_0$ broadly consistent with the historical literature; we flag this result
as a demonstration rather than a validated finding, since the digitized case counts we used have not
been independently verified against the primary source (see \texttt{data/raw/provenance.md}).


\section{Discussion}
The mass-conservation finding (Section~\ref{sec:methods}, Figure~\ref{fig:e02}) is worth dwelling
on because it is a natural but wrong instinct: a modeler who checks $S+I+R \approx N$ as a sanity
test on their SIR solver will see it satisfied to roundoff \emph{even for a badly wrong solve},
because the property follows from the right-hand side summing to zero identically, not from
accurate integration. The phase invariant $Q(t)$ derived from the same exact reduction that gives
our gold standard is the diagnostic that actually carries information, and it is available with no
extra computational cost.

The crossover result (Figure~\ref{fig:e13}) reframes the base paper's implicit assumption. Capaldi
et al.\ do not need to specify a solver precisely because, in the noise regime they study, sampling
uncertainty already dominates; our contribution is to make that dominance relationship explicit and
locate where it stops holding. This is a genuinely two-sided result: it is not simply ``better
solvers are always better'' (uninteresting) or ``solver choice never matters'' (also uninteresting,
and false at low noise) but a quantitative boundary between two regimes, each of which is the
correct practical conclusion in its own part of the parameter space.


\section{Threats to Validity}
\textbf{Profile scale.} The numbers reported here use our \texttt{quick} experiment profile
(Section~\ref{sec:appendix-profiles}), with smaller replicate counts and coarser factor grids than
the full design in our project plan; point estimates of bias and RMSE therefore carry more Monte
Carlo noise than the full run would, though the qualitative orderings (Figure~\ref{fig:e07}'s
convergence-order scaling in particular) are already extremely clean at this scale.

\textbf{Synthetic-data circularity.} Our core claims use data generated from the same SIR model
being fit, which is the correct design for isolating solver effects but cannot by itself show the
model is a good description of any particular real outbreak.

\textbf{Real-data provenance.} The Eyam dataset (E14) was digitized from commonly-reproduced
teaching tables rather than transcribed directly from the primary source in this work; see
\texttt{data/raw/provenance.md}.

\textbf{Single base model family.} While we test SEIR, SIRS, and SIR-with-vaccination (E12), all
four are still relatively low-dimensional, non-stiff systems; whether the crossover phenomenon's
specific numerical values generalize to higher-dimensional or stiffer epidemic models is untested.


\section{Conclusion}
Treating the ODE solver as a first-class experimental variable rather than a black box reveals
structure that a black-box treatment cannot: solver-induced bias in recovered SIR parameters obeys
the solver's own convergence order almost exactly, and this bias competes with noise-induced
variance in a way that can be located precisely as a crossover noise level $\sigma^*$. Below
$\sigma^*$, the numerical-analysis choices this course teaches -- which solver, what step size --
are the dominant factor in whether an epidemiological parameter estimate is trustworthy. Above it,
they are not, and the modeler's effort is better spent on more data or a better noise model.
Reaching this conclusion used the methods of all four CSE 402 labs and the optimization home
assignment (Table~\ref{tab:course}): root finding and LU factorization inside the gold standard
and the implicit solvers, the power method for every eigenvalue and condition number,
least-squares optimization for estimation, quadrature and
interpolation for the reference and observation operators, and random number generation with
Monte Carlo simulation for every noise experiment. Every
figure and number in this report is generated by a single reproducible pipeline
(\texttt{make all}) from the code in this repository.


\bibliographystyle{plain}
\bibliography{refs}

\appendix
\section{Experiment profiles}
\label{sec:appendix-profiles}
Every experiment config (\texttt{configs/*.yaml}) defines a \texttt{quick} and a \texttt{full}
profile. \texttt{quick} is sized to run interactively (seconds to a few minutes per experiment);
\texttt{full} matches the factorial sizes in our original project plan (e.g.\ E07's solver
$\times$ step size $\times$ noise $\times$ sampling interval $\times$ $(\beta,\gamma)$ $\times$
200-replicate design) and is intended for an offline/overnight run. Every number in this report
states which profile produced it.

\section{Reproducibility}
\texttt{make setup \&\& make experiments \&\& make figures \&\& make report} regenerates this PDF
and every figure in it from a clean clone, using the seeds recorded in \texttt{results/manifest.json}.

\section{Team contributions}
\begin{itemize}
  \item \textbf{Saif Uz Zaman} --- numerics core (models, solvers, linear algebra, gold standard).
  \item \textbf{Sayaad Muzahid Masfi} --- report, presentation, and project narrative.
  \item \textbf{Ibtida bin Ahmed} --- estimation and uncertainty quantification.
  \item \textbf{Sakif Naieb Raiyan} --- web dashboard engineering and TypeScript/Python parity.
  \item \textbf{Aurchi Chowdhury} --- experiment design, data provenance, and figure generation.
\end{itemize}


\end{document}
